\section{Introduction}

Stroke is a leading cause of death and of adult disability in the United Kingdom, with more than 100,000 strokes each year. Around 85\% are ischaemic, caused by occlusion of a cerebral artery. Two emergency treatments aim to restore blood flow to the affected brain before the damage becomes permanent: intravenous thrombolysis (IVT), in which a clot-dissolving drug is given by infusion, and mechanical thrombectomy (MT), in which the clot is physically removed through a catheter. The benefit of both depends heavily on how quickly they are delivered.

\subsection{Intravenous thrombolysis}

IVT with alteplase was established by randomised trials showing benefit when given within 3 hours of onset \cite{ninds1995}, later extended to 4.5 hours \cite{hacke2008}. The individual patient data meta-analysis of Emberson et al., pooling 6{,}756 patients from nine trials, quantified how strongly that benefit depends on time \cite{emberson_effect_2014}. The odds of an excellent outcome (modified Rankin Scale, mRS, 0--1 at 3--6 months) were increased by 75\% when treatment was given within 3 hours (OR 1.75, 95\% CI 1.35--2.27) and by 26\% between 3 and 4.5 hours (OR 1.26, 95\% CI 1.05--1.51), with no significant benefit beyond 4.5 hours (OR 1.15, 95\% CI 0.95--1.40). In absolute terms this is roughly ten additional patients per 100 treated achieving an excellent outcome if treated within 3 hours, falling to about five per 100 between 3 and 4.5 hours. Alteplase increased the risk of fatal intracranial haemorrhage in the first week by about 2 percentage points, a risk that did not diminish with later treatment, so the balance of benefit and harm worsens as time passes. IVT can be delivered at any acute stroke unit with access to brain imaging, but it is relatively ineffective at recanalising large, proximal clots.

\subsection{Mechanical thrombectomy}

MT is suitable for patients with large vessel occlusion (LVO) of the anterior circulation, who account for a minority of ischaemic strokes but a disproportionate share of death and severe disability. The HERMES collaboration pooled individual data on 1{,}287 patients from five trials and found that MT substantially reduced disability at 90 days (OR 2.49, 95\% CI 1.76--3.53), with functional independence (mRS 0--2) in 46.0\% of treated patients against 26.5\% of controls, and a number needed to treat of 2.6 for one patient to improve by at least one mRS level \cite{goyal2016}. This makes MT one of the most effective interventions in acute medicine.

As with IVT, the effect diminishes with delay. In the MR CLEAN trial, Fransen et al.\ found that the absolute benefit of MT on functional independence fell by about 6 percentage points for every hour of delay to reperfusion, from an estimated 25.9\% with reperfusion at 3 hours to 6.7\% at 6 hours \cite{fransen_time_2016}. The HERMES time-to-treatment meta-analysis confirmed the pattern across trials: each hour of delay to reperfusion was associated with worse disability (common OR 0.84, 95\% CI 0.76--0.93), and the benefit of MT over IVTas no longer statistically significant beyond 7.3 hours from onset \cite{saver2016}. Later trials showed that patients selected by advanced imaging can still benefit up to 24 hours after they were last known to be well \cite{nogueira2018,albers2018}, but these patients are a selected subgroup, and earlier treatment remains better for them too.

\subsection{The importance of time}

The time dependence of both treatments reflects the underlying biology. In a typical untreated large vessel stroke an estimated 1.9 million neurons are lost each minute \cite{saver2006}. Translated into patient outcomes, each minute saved in delivering IVT has been estimated to give 1.8 extra days of healthy life \cite{meretoja2014}, and each minute saved in delivering MT gives  an extra 4.2 days \cite{meretoja2017}. Time to treatment is therefore a primary determinant of outcome rather than just a secondary operational concern.

\subsection{Implications for the organisation of services}

The two treatments place different demands on a health system. IVT is widely available, whereas MT requires neurointerventional expertise and facilities and is concentrated in a small number of specialist centres; in England it has been provided by 24 neuroscience centres. Patients who live far from one of these centres either travel further to reach it directly, delaying IVT, or attend a local IVT-only unit first for earlier IVT, and are then transferred, delaying MT. Around 10\% of patients admitted with stroke are estimated to be eligible for MT \cite{mcmeekin2017}, and the NHS Long Term Plan set the ambition of reaching this level \cite{nhsltp2019}. In 2024/25, however, 4.4\% of stroke patients received MT and 12.2\% received IVT \cite{ssnap2025}.

\subsection{The addition of two new MT units}

Because the benefit of MT decays so quickly, the number and location of MT units directly affect both how many patients can be treated and how much each of them gains \cite{allen2019}. This report examines the expected effect of adding two new MT units in England, at Taunton and Worcester. In order to isolate the effect of geography, it is assumed all stroke units perform similarly.